% notation.tex -- single source of truth for macros (symbol discipline; define-once).
\newcommand{\R}{\mathbb{R}}
\newcommand{\Z}{\mathbb{Z}}
\newcommand{\N}{\mathbb{N}}
\newcommand{\E}{\mathbb{E}}
\renewcommand{\Pr}{\mathbb{P}}          % probability (\Pr is a built-in operator; renew it)
\newcommand{\Var}{\mathrm{Var}}
\newcommand{\Bin}{\mathrm{Bin}}
\newcommand{\Pois}{\mathrm{Pois}}
\newcommand{\Unif}{\mathrm{Unif}}
\newcommand{\whp}{\textnormal{w.h.p.}}
\newcommand{\eps}{\varepsilon}
\newcommand{\1}{\mathbf{1}}
\newcommand{\supp}{\mathrm{supp}}
\DeclareMathOperator{\polylog}{polylog}

% ---- the model and the two quantities ----
\newcommand{\ondisc}{\mathrm{ondisc}}     % optimal online value
\newcommand{\offdisc}{\mathrm{offdisc}}   % optimal offline value
\newcommand{\colorset}{[k]}
\newcommand{\load}{S}                     % color-class load vector
\newcommand{\rangeop}{R}                  % coordinate pairwise range R_i
\newcommand{\maxmult}{s}                  % |argmax| multiplicity
\newcommand{\rank}{\mathrm{rk}}           % rank = (k-1)(R-1)+s  (was \rho; renamed to avoid clash with signing ratio)

% ---- scales ----
\newcommand{\lam}{\lambda}                % occupancy lambda = Td/(nk) = Theta(d/k)
\newcommand{\Lscale}{L}                   % L = log(ek)
\newcommand{\Hscale}{H}                   % H = loglog n
\newcommand{\Psifn}{\Psi}                 % Psi(lambda,L) = 1 + sqrt(lambda L) + L/log(e+L/lambda)

% ---- shorthands ----
\newcommand{\loglog}{\log\log}
\newcommand{\Oh}{O}
\newcommand{\Om}{\Omega}
\newcommand{\Th}{\Theta}
